\cite{negri2017review} - digital teins moniyret got anomalies, geomnetrics and plastic deformation on matieral of the physical twin, model reliabilkity, study long term behaviour of the sytstem and perfict performance, proivide infgormation across liscycle phases, optimize system behaviour during the design phase, codnition based maintenance, diagnostics, prognostic activites, prediction of failures of the systemn, 

\cite{tao2018digital} - DTs for modelling and simualtion,d ata fusuion from processing massive volumes of data collectied from machines and the environmenty, so mainly for data generation, modelling, cleaning, clustering, mining, and evolution. 


\textcite{ghanbarifard2023digital} further reinforce this perspective by discussing distributed DT composition in dynamic, evolving operational contexts. ntegrated DTs in complex operations environments
With the advent of complex operations environments, a
discussion is necessary concerning the use of distributed DT,
considering how inevitable it is. Thus, a comprehensive DT is
incumbent upon the result of integrating all DT parts.
A DT includes all relevant information of the system and is
comprised of interconnected digital artifacts that are linked
together. An intertwined model can merge new behaviors than
partial models. In complex systems, integrating all parts models
is valuable as it allows for the study of complex behaviors that
arise from the exchange of information or influences among
different components of the system [65, 66]. This integration
facilitates a more comprehensive understanding of the system as
a whole.
Vrabic et al. [66] proposed that partial models of a DT can
be merged through a unified model space that defines their
interrelations. They found that this integration enhances the
value of the DT at various Points along the lifecycle, including
the assessment of efficiency during the design and monitoring
of the physical infrastructure in real time.
To ensure that different components from various value
chains can work together in Industry 4.0, it needs to incorporate
DTs into the SOS. To achieve the necessary flexibility, these
digital representations must be seamlessly integrated into their
respective environments. Gollner et al. [33] recognized the
importance of a holistic Systems Engineering approach for
designing an architecture for interoperable DTs. To facilitate the
deployment of compatible DTs in Industry 4.0, they developed
guidelines and a modeling approach.
In the context of a shop floor, Liu et al. [67] proposed an
approach for constructing a comprehensive DT model, which is
a complex system. This method involves the use of different
models containing system language, the "MagicGrid" [68], and
the "V model" [69] to gradually construct an integrated DT for
the shop floor.
Zohdi [70] proposed a machine-learning framework that
utilizes DT to rapidly simulate fire propagation in complex
environments. This framework demonstrates the integration of
machine-learning algorithms with multiple-step fire spread
models The DT can continuously update itself in near real-time
due to a simulation technique that can be rapidly linked to
telecommunications, cameras, and sensors. It is now possible to
drive fully autonomously on a single machine as SoS [71].
Bertoni et al. [72] introduced a technique for merging systems
and SoS simulations using DTs to compare different product-
service systems configurations and integrate variables into
simulations to build a DT.
Despite the use of project management techniques, cost
enhancement, and schedule delays often persist for industries
managing complex projects [73]. These disciplines can
frequently operate independently, causing discrepancies in
product features, keeping track of design alterations, and
decision-making obstacles. Jonkers et al. [74] proposed a
management flight simulator using Dt that combines sub-
models from various disciplines and considers the cultural
factors that impact design project management and change
decision



\cite{vergara2023federated} - Nevertheless, much research is needed to develop and validate
appropriate mechanisms to ensure its successful deployment
in complex real-world cases that require collaboration among
individual systems. A Federated Digital Twin (FDT) has been
identified as a promising solution for this approach, since it
allows the interconnection among autonomous DTs in the virtual
space, leveraging their advantages and enabling interaction,
collaboration and shared learning. Additionally, since a FDT is
envisaged as a network of cooperative DTs, cognitive principles
can be applied to assist the overall operations through knowledge
acquisition and reasoning, leading to an informed and intelligent
decision making
A Federated Digital Twin (FDT) consists of the integration
of autonomous DTs interconnected in a virtual environment,
governed by a set of rules that determine interoperability,
coordination and secure communications. An FDT provides a
holistic view of a complex system and a broader understanding
of the relationship between the virtual entities (DTs) based on
aspects such as interaction, collaboration and synchronization.
The federation generates a virtual medium for data ex-
change between the DTs, which enables cooperation to achieve
common goals, according to the individual and collective
contributions. Furthermore, it can be used to perform sim-
ulations, predictions and data analytics to optimize system
operations and improve decision making in the federation,
before transferring it to the physical entities. Since each DT
has been designed and implemented with specific objectives
and supervises its own monitoring and virtual representation,
the federation enhances decentralization, scalability and share
B. Architectural styles
• Centralized: The most basic FDT architecture, managed
by a central orchestrator, which supervise the communi-
cations, coordination, synchronization and goals.
• Peer-to-Peer: The federation is managed without a central
orchestrator, instead, the DTs interact directly in a P2P
fashion. They should self-coordinate and self-synchronize
to achieve common federation goals, in a more sophisti-
cated architecture.
• Hierarchical: Large setting of an FDT, where the oper-
ations, coordination, data exchange and synchronization
are managed by a global central orchestrator, and where
local federations supervise operations within their scope.
• Regional: Large setting of an FDT where the operations,
coordination, data exchange and synchronization should
be managed between local federation managers (in a P2P
fashion), since there is no a global central orchestrator


\cite{pias2025scaling} -  First, we summarise increased expectations as a hierarchy of scales (1–4):
1. System-specific: Typical procedures and tools for the monitoring and simulation of a single system
or asset, relating to conventional asset management, prognostics, and control (Bruynseels et al.,
2018; Barricelli et al., 2019; Furini et al., 2022). The boundaries of the virtual representation (or the
limits of the DT) are intuitive, with clear approximations.
2. Local collection: Aligned with the ideas from the Internet of Things, including connected,
distributed representations on a smaller scale. Examples include industrial applications (Kappel
et al., 2022), urban farming in the United Kingdom and the United States (Jans-Singh et al., 2020),
and land-based aquaculture (Lima et al., 2022). The boundary of the representation becomes less
clear—for example, in agriculture, the soil depth and surrounding environment become increas-
ingly complex.
3. Ecosystems: Typical applications include urban analytics, smart infrastructure (Oughton and
Russell, 2020), population-based structural health monitoring (Worden et al., 2020), and Industry
5.0 applications (Braque et al., 2021). More recently, large-scale applications in health, medicine,
and education are emerging (Björnsson et al., 2020; Furini et al., 2022). While these aggregates are
typically constituted of engineered systems, with defined boundaries, the limit on the extent of the
DT representation is harder to define.
4. Planetary: Distributed and parallel DTs combined to accomplish scaled-up tasks of complex
interconnected systems (sensing for weather, pollution, large-scale energy modelling, etc.). Highly
complex boundaries and interactions are associated with the aggregat
3.2. An example of increasing scale: wind energy
Figure 2 presents one natural example of DT aggregation, progressing through the hierarchies of scale
with lighter shades of purple.
1. System-specific: A single turbine—even at this scale, the DT is an aggregate of data + models, for
example, corresponding to different components and their interactions (the nacelle, blades, and
foundations).
2. Local collection: The wind farm of constituent turbines. The boundary becomes more complex:
where is the limit of representation and how do we capture the interactions between other turbines,
the grid, and the environment (on or off-shore)?
3. Ecosystem: Alternative power sources are now included, alongside human agents. Interoperability
of DT information becomes key, especially between organisations. The boundary extends to both
natural (human as an agent) and engineered system behaviour.
4. Planetary: Repeated extensions (into the background). The DT combines several whole energy
systems; this could be international. Note that the elements of the aggregate can be shared between
parts (e.g., the weather models)



\cite{kunnari2025distributed} - When multiple DTs work on the same physical entity, they need to communicate and coordinate with
each other as to not cause problems within the physical system. Without a structured coordination
mechanism, DTs risk conflicting optimization goals, conflicting commands and inconsistent real-time
behavior. To address these challenges, a hierarchical structure becomes necessary, leading to the
Distributed, Hierarchical Digital Twin (DHDT) paradigm. A DHDT could take three forms: (1) a
comprehensive DT that integrates multiple specialized DTs into a cohesive system, (2) a network of
interconnected DTs that collectively represent a complex production process, or (3) a mix of the first two
approaches. In all cases, a hierarchical structure emerges where a supervisory DT governs optimization
goals while subordinate DTs contribute domain-specific insights. DHDTs offer a more granular
representation of production processes, leading to potentially better results.
However, the integration of multiple DTs into a DHDT potentially extends challenges encountered in
traditional DT implementations. While existing research has addressed these challenges in the context of
standalone DTs, they require renewed examination within the DHDT paradigm. Key challenges include
the integration of new components, synchronization of distributed simulation models, communication
between digital and physical assets, scalability, and distributed governance. The interplay between
multiple DTs introduces complexities in interoperability, real-time coordination, and cross-domain
decision-making that are not present in isolated DT environments.
3.1 Distributed Digital Twins
A targeted search within major academic databases was performed using key terms such as "distributed
digital twins," "federated digital twins," and "multi-agent digital twins." This process yielded 24 potential
sources, with 16 ultimately included in the final review. To extract patterns across these publications, an
initial concept matrix was developed and later transformed into a morphological box to support a more
structured representation of recurring themes. The dimensions of the morphological box were not selected
arbitrarily; rather, they emerged inductively during the early analysis phase. Recurring topics such as
1508
Authorized licensed use limited to: McMaster University. Downloaded on March 07,2026 at 16:07:07 UTC from IEEE Xplore. Restrictions apply.
Kunnari and Strassburger
governance models, architectural types, and optimization objectives were consistently observed across the
sources. Their frequency and relevance prompted the structured cross-comparison that underpins Table 1.
The left column of the morphological box shows six different categories that were analyzed in the
reviewed literature. Many papers fell into multiple varieties within a single category. The numbers listed
in the table indicate how many sources addressed each topic. The categories are:

\cite{yu2025federated} - Urban logistics faces increasing pressure from rising population densities, escalating delivery demands, and
constrained urban resources. Traditional logistics systems struggle to adapt to real-time urban dynamics, leading
to inefficiencies, congestion, and environmental concerns. A key challenge lies in mobilizing underutilized assets,
such as off-hour freight parking, and adopting multimodal solutions to navigate diverse and increasingly strict
regulations, thereby enhancing both sustainability and operational efficiency. However, effective management
and utilization of these assets require real-time visibility, cross-stakeholder collaboration, and intelligent
decision-making. This study proposes a federated digital twin platform to enhance logistics operations efficiency
by integrating asset management and knowledge-driven operations management, relying on real-time asset
visibility and delivery knowledge, such as destination characteristics and preferred logistics modalities. Unlike
traditional logistics planning, which relies on static assumptions, our approach adapts to urban constraints by
continuously querying real-time asset information and integrating logistics-related knowledge into operations
management. To assess the effectiveness of this approach, an optimization-based simulation framework with
decision-making tools is developed. The study evaluates multi-echelon logistics networks, incorporating micro-
hubs, dynamic transshipment points, and multimodal logistics options, including on-foot porters, E-cargo bikes,
and Road Autonomous Delivery Robots (RADRs). Findings demonstrate that integrating federated digital twins
with knowledge-driven approaches, such as destination-based clustering and modality selection, reduces costs by
over 50 % and emissions by more than 30 %. This study underscores the transformative potential of digital twins
in enabling real-time, knowledge-driven operations management, and fostering more sustainable and efficient
urban logistics systems




% ===============================================
\cite{alma991032896202207371} - Reliability: Reliability is the probability that a product will operate or a service will be
provided properly for a specified period of time (design life) under the design operating
conditions (such as temperature, load, volt…) without failure.
In other words, reliability may be used as a measure of the system’s success in pro-
viding its function properly during its design life. 


\cite{ram2013system} - It is well known that reliability is the probability of a
device performing its purpose adequately for the period
intended under the given operating conditions. Reliability
expressed as a probability: probability is the ratio of the
number of times. We can expect an event to occur to the
total number of trials undertaken. The maximum value of
this fraction is one and minimum value is zero. A proba-
bility of one means a certainty i.e. the expected event
occurs in almost every case. The expected event is the
adequate or satisfactory performance of the device. The
reliability factor of one means that the device performs
satisfactory for the prescribed duration under the given
environmental conditions. Similarly, a reliability factor of
zero means that in almost in all cases the equipment would
fail to meet the required performance level. Adequate
performance means what is expected of a device or system
duration of adequate performance represents a measure of
the period for which the performance is satisfactory. The
operating conditions in which we expect a device to
function adequately could be with regard to temperature,
humidity, shock, vibration, etc. Reliability can be extended
to include several levels of system performance. The first
performance-oriented extension of reliability is to replace a
single exportable-level of operation by a set of perfor-
mance level. This approach is used for evaluating network
performance and reliability. The performance level is
based on matrix derived from an application-dependent
performance model. The performance level might be the
rate of job completion, the response time or the number of
jobs completed in a given time interval. Under a wide
range of conditions, the consequences of failure are grave,
but there is a correspondingly higher budget. A pencil
sharpener may be more reliable than an airliner but has
many different sets of operating conditions, insignificant
consequences of failure and a much lower budget. A reli-
ability program is used to document exactly what best
practices, close tasks, methods, tools, analysis & tests are
required for a particular system as well as to clarify cus-
tomer requirements for reliability assessment. Reliability
requirements or specifications, test plan, and contract
statements of the system/subsystem. Maintainability
requirements address system issue of costs as well as time
to repair. Reliability prediction is the combination of the
creation of the proper reliability model together with esti-
mating the input parameters (like failure rates for a par-
ticular failure mode or event and the MTTR the system for
a particular failure) for this model and finally to provide a
system-level estimate for output reliability parameters
(system availability or a particular function failure fre-
quency). Requirements are specified using reliability
parameters. The most common parameter of reliability is
MTTF. This parameter is very useful for systems that are
operated frequently such as vehicles, machinery & elec-
tronic equipments. Reliability increases as the MTTF
increases. Reliability modeling is the process of predicting
or understanding the reliability of a component or system
prior to its implementation. Two types of analysis are often
used to model a system’s reliability behavior are fault tree
analysis and reliability block diagrams. Reliability must be
‘‘designed in’’ to the system. During system design, the
top-level reliability requirements are then allocated to
subsystems by design engineers, maintainers and reliability
engineers working together. Reliability design begins with
the development of a (system) model. Reliability model
use block diagrams and fault trees to provide a graphical
means of evaluating the relationship between different
parts of the system (Verma et al. 2010).
Wang and Pham (1997) analyzed and summarized a
survey with the typical existing Monte Carlo reliability,
availability and MTBF simulation procedures, variance
reduction methods, and random variate generation algo-
rithms. The very important and extensive explanation was
discussed by Avizienis et al. (2004) in the form of basic
concepts and taxonomy of dependability and secure com-
puting. Here author tried to give the foremost fields and
approaches of engineering and sciences, alphabetically
where the reliability theory has been applied.




\cite{laprie1996software} - 
impairements to dependability are faults, errors, and failrture. Methods to ensure dependability are through fasult prevention, removal, tolerance, and forcasting. A system transitions between correct service and incorrect service, relaibility is a measure of the contitnious delievry of the correct service.